\documentclass[12pt]{article}
\usepackage[utf8]{inputenc}
\usepackage{lipsum} % to generate some filler text
\usepackage{fullpage}
\usepackage{bm}
\usepackage{amsmath,amssymb,dsfont, mathtools}
\usepackage[pdftex]{graphicx}
\usepackage{epstopdf}
\usepackage{subcaption}
\usepackage{natbib}
\usepackage{hyperref}
\usepackage{placeins}
\usepackage{algorithm}
\usepackage{algorithmic}
\usepackage{amsmath}
\usepackage{booktabs}


% import Eq and Section references from the main manuscript where needed
% \usepackage{xr}
% \externaldocument{manuscript}
\DeclareMathOperator{\plim}{plim}
\DeclareMathOperator{\vect}{vec}
\DeclareMathOperator{\D}{D}
\DeclareMathOperator{\Diag}{Diag}
\DeclareMathOperator{\argmin}{argmin} 
\DeclareMathOperator{\argmax}{argmax} 
\newcommand{\bs}{\boldsymbol}
\newcommand{\E}{\mathbb{E}}
\newcommand{\tr}{\mathrm{tr}}
\newcommand{\abs}{\mathrm{abs}}
\newcommand{\var}{\mathrm{Var}}
\newcommand{\mb}{\mathbb}
\newcommand{\mf}{\mathbf}
\newcommand{\mc}{\mathcal}
\newcommand{\mk}{\mathfrak}
\renewcommand{\vec}{\text{vec}}
\newcommand{\red}[1]{{\color{red}#1}} 
\newcommand{\h}{\mathbf{h}}
\newcommand{\Hh}{\mathbf{H}}
\newcommand{\W}{\mathbf{W}}
\newcommand{\M}{\mathbf{M}}
\newcommand{\R}{\mathbf{R}}
\newcommand{\U}{\mathbf{U}}
\newcommand{\V}{\mathbf{V}}
\newcommand{\I}{\mathbf{I}}
\newcommand{\Y}{\mathbf{Y}}
\newcommand{\X}{\mathbf{X}}
\newcommand{\Ss}{\mathbf{S}}
\newcommand{\A}{\mathbf{A}}
\newcommand{\J}{\mathbf{J}}
\newcommand{\K}{\mathbf{K}}
\newcommand{\Pp}{\mathbf{P}}
\newcommand{\0}{\mathbf{0}}
\newcommand{\C}{\mathbf{c}}
\newcommand{\f}{\mathbf{f}}
\newcommand{\Lag}{\mathbf{L}}
\newcommand{\B}{\mathbf{B}}
\newcommand{\s}{\mathbf{s}}
\newcommand{\Z}{\mathbf{Z}}
%Bold greek letter using the first two letters
\newcommand{\be}{\boldsymbol{\beta}}
\newcommand{\la}{\boldsymbol{\lambda}}
\newcommand{\rh}{\boldsymbol{\rho}}
\newcommand{\et}{\boldsymbol{\eta}}
\newcommand{\mm}{\boldsymbol{\mu}}
\newcommand{\de}{\boldsymbol{\delta}}
\newcommand{\xx}{\boldsymbol{\xi}}
\newcommand{\bb}{\boldsymbol{\beta}}
\newcommand{\Om}{\boldsymbol{\Omega}}
\newcommand{\Si}{\boldsymbol{\Sigma}}
\newcommand{\Up}{\boldsymbol{\Upsilon}}
\newcommand{\Th}{\boldsymbol{\theta}}
\newcommand{\Ps}{\boldsymbol{\Psi}}
\newcommand{\Ph}{\boldsymbol{\Phi}}
\newcommand{\La}{\boldsymbol{\Lambda}}
\newtheorem{rmk}{Remark}

\newcommand{\xvec}{\boldsymbol}
\newcommand{\xmat}{\mathbf}
\newcommand{\xset}{\mathds}

\newcommand{\y}{\xvec{Y}}
%\newcommand{\K}{\mathbf{K}}
\newcommand{\x}{\xvec{x}}
\newcommand{\z}{\xvec{z}}
\DeclareMathOperator{\e}{\varepsilon}
\usepackage[shortlabels]{enumitem}
\renewcommand{\theenumi}{\Alph{enumi}}

\newcommand{\blue}{\textcolor{blue}}

% package needed for optional arguments
\usepackage{xifthen}
\usepackage{float}
% define counters for reviewers and their points
\newcounter{reviewer}
\setcounter{reviewer}{0}
\newcounter{point}[reviewer]
\setcounter{point}{0}

% This refines the format of how the reviewer/point reference will appear.
\renewcommand{\thepoint}{P\,\thereviewer.\arabic{point}} 

% command declarations for reviewer points and our responses
\newcommand{\reviewersection}{\stepcounter{reviewer} \bigskip \hrule
                  \section*{}}

\newenvironment{point}
   {\refstepcounter{point} \bigskip \noindent {\textbf{Point~\thepoint} } ---\ }
   {\par }

\newcommand{\shortpoint}[1]{\refstepcounter{point}  \bigskip \noindent 
	{\textbf{Reviewer~Point~\thepoint} } ---~#1\par }

\newenvironment{reply}
   {\medskip \noindent \begin{sf}\color{Blue}\textbf{Reply}:\  }
   {\medskip \end{sf}}

\newenvironment{replycont}
   {\medskip \noindent \begin{sf}\color{Blue}\  }
   {\medskip \end{sf}}


 \newenvironment{changes}
   {\medskip \noindent \begin{sf}\color{ForestGreen}\textbf{Changes}:\  }
   {\medskip \end{sf}}

\newcommand{\shortreply}[2][]{\medskip \noindent \begin{sf}\textbf{Reply}:\  #2
	\ifthenelse{\equal{#1}{}}{}{ \hfill \footnotesize (#1)}%
	\medskip \end{sf}}

\usepackage[pagewise, mathlines]{lineno} % Needed to include, otherwise, doesn't work
\usepackage[dvipsnames]{xcolor}

\begin{document}

\author{ }
\bibliographystyle{agsm}
\date{}
\title{Responses to reviewers}
\maketitle

\vspace*{-1.0cm}

\begin{sf}
\noindent \noindent We thank the reviewer for the insightful suggestions provided during this revision round. We have addressed each of the points raised; a detailed point-by-point response is provided below. 

\noindent The revised manuscript incorporates several major enhancements to address the concerns raised and to improve the overall transparency of the approach:
\end{sf}

\reviewersection

\begin{sf}
\noindent We sincerely thank the reviewer for the careful second-round evaluation and for the positive assessment of the revised manuscript. We are grateful for the acknowledgment that the manuscript has substantially improved and that the methodological concerns raised in the previous round have been satisfactorily addressed.

\noindent We especially appreciate the reviewer’s recognition of the new validation experiment, the ordering/sparsity analysis, the clarification of the covariance structure, the expanded baseline comparisons, and the additional sensitivity analyses. We also thank the reviewer for the remaining constructive suggestions regarding reproducibility, presentation, and proofreading. 

\noindent Below we provide a detailed point-by-point response to all remaining comments. All corresponding modifications have been incorporated into the revised manuscript.
\end{sf}

%%%%%%%%%%%%%%%%%%%%%%%%%%%%%%%%%%%%%%%%%%%%%%%%%%
%%%%%%%%%%%%%%%%%%%% POINT 1 %%%%%%%%%%%%%%%%%%%%%
%%%%%%%%%%%%%%%%%%%%%%%%%%%%%%%%%%%%%%%%%%%%%%%%%%

\begin{point}
In the response letter, the authors now provide a clear description of the optimization routine, including BFGS via fminunc, log-domain parameterization of the hyperparameters, multi-start initialization, sequential initialization from simpler models, and the convergence tolerances used. Since these are exactly the computational details requested in the previous round, it would be helpful to include a concise version of them in the manuscript itself, either in Section 2 or in an appendix. If an optimization summary can be adapted into the paper, that would make the implementation substantially easier to reproduce.
\end{point}


\begin{reply}
We thank the reviewer for this helpful suggestion. We agree that including a concise description of the optimization strategy directly in the manuscript improves reproducibility and implementation transparency. In the revised version, we have added a new \emph{Optimization Summary} subsection in the appendix (reproduced below for the reviewer's convenience). The same information, together with the executable code, is also available in the public repository accompanying the paper.
\end{reply}

\begin{changes}
The following subsection has been added to the appendix.

\medskip
\noindent\textbf{Appendix~X. Optimization summary.}
All hyperparameters of the MFGP model are estimated by maximum likelihood (or restricted maximum likelihood, where indicated) using the BFGS quasi-Newton solver provided by MATLAB's \texttt{fminunc} (Optimization Toolbox). The objective is the (Vecchia-approximated) negative log marginal likelihood; gradients are computed by central finite differences. The implementation details that are common to all experiments are summarized below.

\smallskip
\noindent\textit{Parameterization.} The hyperparameter vector $\bs{\theta}\in\mathbb{R}^{p}$ ($p=11$ for the stationary configurations and $p=14$ when the discrepancy length-scale is modeled by the auxiliary GP discussed in Section~2) collects all variance, length-scale, and noise parameters. To enforce positivity and to keep the optimization landscape on a comparable scale, the optimization is carried out in the log-domain: each positive parameter $\phi>0$ is reparameterized as $\theta=\log\phi$ and recovered as $\phi=\exp(\theta)$ inside the likelihood (see, e.g., the variance and length-scale terms \mbox{\texttt{sigma\_LF\_t = exp(hyp(1))}}, \mbox{\texttt{ell\_LF\_t = exp(hyp(2))}}, etc. in the public code).

\smallskip
\noindent\textit{Solver and options.} We use the quasi-Newton (BFGS) algorithm of \texttt{fminunc} with the following options (MATLAB syntax):

\begin{center}
\small
\begin{tabular}{ll}
\toprule
Option & Value \\
\midrule
\texttt{Algorithm}                   & \texttt{'quasi-newton'} (BFGS) \\
\texttt{SpecifyObjectiveGradient}    & \texttt{false} \\
\texttt{FiniteDifferenceType}        & \texttt{'central'} \\
\texttt{FiniteDifferenceStepSize}    & $10^{-4}$ \\
\texttt{FunctionTolerance}           & $10^{-8}$ \\
\texttt{StepTolerance}               & $10^{-8}$ \\
\texttt{MaxIterations}               & 200 (synthetic) / 50 (real data) \\
\texttt{MaxFunctionEvaluations}      & 5000 \\
\texttt{TypicalX}                    & $\mathbf{1}+|\bs{\theta}^{(0)}|$ (componentwise) \\
\bottomrule
\end{tabular}
\end{center}

The smaller \texttt{MaxIterations} budget on the real-data experiment is paired with the sequential-initialization strategy described below, which restarts each fit from a previously converged set of hyperparameters and therefore requires only a few BFGS updates to refine.

\smallskip
\noindent\textit{Initialization and multi-start.} For the first run on a given station/configuration, the starting point $\bs{\theta}^{(0)}$ is drawn uniformly from $[0,1]^{p}$ (in the log-domain). To reduce sensitivity to the initial guess on the synthetic experiments, the fit is repeated from $K$ independent random starts (we use $K=5$ in Tables~1--2; results are robust to $K\in\{3,\dots,10\}$) and the minimizer with the smallest final objective is retained.

\smallskip
\noindent\textit{Sequential initialization from simpler models.} On the real-data benchmark we fit the six MFGP variants in increasing order of complexity (constant\,$\rho$ $\rightarrow$ adaptive\,$\rho$ $\rightarrow$ warped variants $\rightarrow$ GP-driven $\rho(s)$). The optimum $\hat{\bs{\theta}}$ of each simpler model is used as the starting point of the next, more complex one (the extra hyperparameters introduced by the more flexible model are initialized at small random values). The same warmstart is also used \emph{across} the leave-one-station-out folds, where the previously fitted hyperparameters of the same configuration on the other 17 stations serve as $\bs{\theta}^{(0)}$ (see \texttt{ResultsHistory} mechanism in the public code).

\smallskip
\noindent\textit{Convergence diagnostics.} Convergence is declared when both \texttt{FunctionTolerance} and \texttt{StepTolerance} are satisfied, or when the first-order optimality measure falls below $10^{-6}$. Runs that terminate at \texttt{MaxIterations} without satisfying these criteria (less than $5\%$ of all fits in our experiments) are flagged and re-run from a different random start.

\smallskip
\noindent\textit{Software.} All optimization routines run under MATLAB R2023a, requiring only the Optimization Toolbox (\texttt{fminunc}) and the Statistics and Machine Learning Toolbox (\texttt{fitrgp}, used for the single-fidelity GP baseline).
\end{changes}

%%%%%%%%%%%%%%%%%%%%%%%%%%%%%%%%%%%%%%%%%%%%%%%%%%
%%%%%%%%%%%%%%%%%%%% POINT 2 %%%%%%%%%%%%%%%%%%%%%
%%%%%%%%%%%%%%%%%%%%%%%%%%%%%%%%%%%%%%%%%%%%%%%%%%

\begin{point}
The response letter also provides a useful quantitative summary of negative predictions across the fitted models, showing that 90 of the 108 fitted models produced no negative values, that most of the remaining cases were of very small magnitude, and that only one adaptive-$\rho(s)$ configuration (Station 856) produced a more substantial negative excursion ($-6.33$). I think some brief version of this diagnostic summary would be worth including in the manuscript or appendix, since it helps justify the discussion of warping and clarifies how often physically implausible predictions arise in practice.
\end{point}

\begin{reply}
We thank the reviewer for this valuable suggestion. We agree that including a concise diagnostic summary on negative predictions improves the practical interpretation of the methodology and clarifies how often physically implausible predictions actually arise. Accordingly, we have added a dedicated appendix (\emph{Appendix: Empirical diagnostic of negative predictions}, label \texttt{app:negatives}) reporting Table~\ref{tab:detailed_negatives} together with a short discussion: $90/108$ fits ($83.3\%$) produce no negative values, an additional $14$ fits ($13.0\%$) yield $1$--$3$ negatives of average magnitude $\sim 10^{-1}$, and only one adaptive-$\rho(\bm{s})$ configuration (Station~856) produces a substantial outlier at $-6.33$. A pointer to this appendix has also been inserted in the real-data results discussion (Section~\ref{Experiments}), so that the reader encounters the diagnostic directly after the main predictive-accuracy comparison.
\end{reply}

\begin{changes}
The following appendix has been added to the manuscript (verbatim, abridged here for the response letter).

\medskip
\noindent\textbf{Appendix~Y. Empirical diagnostic of negative predictions.}
Across the $108$ fits of the leave-one-station-out experiment ($6$ MFGP variants $\times$ $18$ stations), $90$ fits ($83.3\%$) produce no negative predicted values, $14$ ($13.0\%$) produce $1$--$3$ small-magnitude negatives (average minima between $-0.08$ and $-0.22$), and only $4$ fits exhibit larger excursions, dominated by a single outlier: the adaptive-$\rho(\bm{s})$ MFGP configuration with constant GLS mean on Station~856, which reaches a minimum of $-6.33$. The breakdown by number of negatives per fit is reported in the following table.

\begin{center}
\small
\begin{tabular}{@{}lcc@{}}
\toprule
Error Category & No.\ of Models & Average Minimum Value \\ \midrule
Models with 0 negatives    & 90 & \phantom{$-$}0.0000 \\
Models with 1 negative     & 8  & $-0.0818$ \\
Models with 2 negatives    & 3  & $-0.2232$ \\
Models with 3 negatives    & 3  & $-0.0812$ \\
Models with 4 negatives    & 1  & $-0.1597$ \\
Models with 5 negatives    & 1  & $-0.4586$ \\
Models with 6 negatives    & 1  & $-0.5877$ \\
Models with $>6$ negatives & 1  & $-6.3333$ \\ \bottomrule
\end{tabular}
\end{center}

The appendix concludes that physically implausible predictions are infrequent ($\approx 17\%$ of fits) and, when they occur, of negligible magnitude in all but one configuration, so that the optional warping post-processing step is sufficient to enforce positivity in practice without modifying the core algorithm.
\end{changes}

%%%%%%%%%%%%%%%%%%%%%%%%%%%%%%%%%%%%%%%%%%%%%%%%%%
%%%%%%%%%%%%%%%%%%%% POINT 3 %%%%%%%%%%%%%%%%%%%%%
%%%%%%%%%%%%%%%%%%%%%%%%%%%%%%%%%%%%%%%%%%%%%%%%%%

\begin{point}
For readability, it may be worth aligning the placement of Table B.2 with the appendix in which it is mainly discussed. At present, Table B.2 is discussed in Appendix B.1 but placed in Appendix F, which creates a bit of back-and-forth for the reader.
\end{point}

\begin{reply}
We thank the reviewer for pointing this out. The cross-reference inconsistency was caused by two issues in the LaTeX source: (i) the \texttt{\textbackslash label} of the ordering-strategy table was placed \emph{outside} the \texttt{table} environment, so the table was numbered as part of Appendix~F instead of Appendix~B; and (ii) the float placement specifier \texttt{[ht]} allowed the table to drift away from its discussion. We have now (a) moved the label inside the table (as \texttt{tab:OrderingStrategies}), (b) tightened the float placement to \texttt{[!htbp]} and added a \texttt{\textbackslash FloatBarrier} so that the table is rendered within Appendix~\ref{sec:supp_ordering_accuracy} where it is discussed, and (c) updated the introductory paragraph of the \emph{Additional results} appendix and the in-text reference in Appendix~B accordingly. Table B.2 now appears in the same appendix as its discussion, removing the back-and-forth highlighted by the reviewer.
\end{reply}

%%%%%%%%%%%%%%%%%%%%%%%%%%%%%%%%%%%%%%%%%%%%%%%%%%
%%%%%%%%%%%%%%%%%%%% POINT 4 %%%%%%%%%%%%%%%%%%%%%
%%%%%%%%%%%%%%%%%%%%%%%%%%%%%%%%%%%%%%%%%%%%%%%%%%

\begin{point}
The Figure 1 caption could be slightly more self-contained by clarifying what the shaded envelopes represent, for example whether they correspond to predictive intervals, pointwise uncertainty bands, or something else.
\end{point}

\begin{reply}
We thank the reviewer for this suggestion. We have revised the caption of Figure 1 to clarify the interpretation of the shaded envelopes and to make the figure more self-contained for readers.
\end{reply}

%%%%%%%%%%%%%%%%%%%%%%%%%%%%%%%%%%%%%%%%%%%%%%%%%%
%%%%%%%%%%%%%%%%%%%% POINT 5 %%%%%%%%%%%%%%%%%%%%%
%%%%%%%%%%%%%%%%%%%%%%%%%%%%%%%%%%%%%%%%%%%%%%%%%%

\begin{point}
A final proofreading pass would still be beneficial. A few issues I noticed include, non-exhaustively:
\begin{itemize}
    \item “We iused the following metrics” (Appendix B.1).
    \item “coeefficient” in the Figure 2 caption.
    \item Algorithm 1, line 17: duplicated “if if GLS == False then”.
    \item Algorithm 1 appears to be missing matching end if/end statements for the inner conditional and the outer block.
\end{itemize}
These are all minor and should be easy to fix in a final polish.
\end{point}

\begin{reply}
We thank the reviewer for the careful proofreading and for identifying these issues. We have corrected all typographical and formatting errors mentioned above, including the wording issues, the duplicated conditional statement in Algorithm 1, and the missing matching end statements. We also performed an additional proofreading pass throughout the manuscript to improve overall clarity and consistency.
\end{reply}





\end{document}